\chapter{Open questions}\label{ch:not-known-index}

Open items about submission below Metal stay in \cref{ch:submission-below-metal}, where they arise, except where they
bear on residency or ordering.

\section{Device, registers and encoding}\label{sec:device-registers-encoding}

\begin{itemize}
\item Whether a released register can read its old value; no sentinel-checked probe of the old-value branch has been run (MM~\mm{25.145}).
\item What IR0-IR15 hold and whether any executed program names one (Apple's class table; MM~\mm{25.59}).
\item Whether R126/R127 are reserved or the part has only 126 of 128 registers (MM~\mm{24.2} row H1).
\item Physical register capacity and allocation granularity per core (MM~\mm{24.2} row H4, MM~\mm{20}, MM~\mm{25.20}, MM~\mm{25.49}).
\item Why this compiler's residency counter reads about twice Apple's kernel's at high register counts (64–69 against 32–45 per core) (MM~\mm{25.121}).
\item Simultaneous residency count for 1,024-thread threadgroups: the forward-progress probe's (\code{fp.m}) 200 GPU-wide,
  which implies at least 10 per core, against the driver's 60 GPU-wide, 3 per core (MM~\mm{25.141.1}, MM~\mm{25.141.20}; \cref{ch:device-execution-model}).
\item Whether the encoding holes of \op{1004} and \op{798} are unreachable in the encoding or only in the field maps (MM~\mm{25.133}, MM~\mm{25.196}).
\item Release order inside instructions other than the measured ALU forms and the word store (MM~\mm{25.117}, MM~\mm{25.145}).
\item Which narrow-pool policy gives better occupancy (MM~\mm{25.146}).
\item Which lengths class-7 instructions take when byte6 = 0xa1 (MM~\mm{25}); which field selects width 10 in \op{11842}; a sufficient condition for width 14 in \op{17235}/\op{17244} (MM~\mm{25.62}).
\item \op{12674} residue bit byte1{[}6{]} beyond the rendering split (MM~\mm{25.32}).
\item Bits 33 and 37 of operand 1 in ALU and load forms (the running example in \cref{sec:field-maps}).
\item Whether an independent decoder can be derived from the per-length field maps, with encode-then-compare refusal (MM~\mm{25.209}).
\item Whether a wrong slot-32 value with correct presence affects any kernel shape beyond those tested (MM~\mm{10}).
\item The 12 open class-constant bytes (0xC0000000 at 1304, word 1 at 1324, 0x40000004 at 1440) (MM~\mm{25.148}).
\item Q in the scalar metadata class, the per-kernel multiple of 4 that slot 4 carries as Q $-$ 4 and that correlates with no readable field (\code{agxforge/g17/mdgen.py}).
\item The main-program rule behind the 70 L3 misses (MM~\mm{25.120}, MM~\mm{25.120.1}).
\item Whether a non-32-bit uniform write occupies one uniform (MM~\mm{25.120}, MM~\mm{25.120.1}).
\item Why 442 constant-program uniform writes are never read (MM~\mm{25.120}, MM~\mm{25.120.1}).
\item Whether the register count in the below-Metal USC Registers packet affects residency (MM~\mm{25.143}; \cref{ch:submission-below-metal}).
\item Slot 44's byte position rule across classes (MM~\mm{25.89}, MM~\mm{25.120}, \code{agxforge/g17/mdgen.py}).
\end{itemize}

\section{Scalar code and cross-lane operations}\label{sec:scalar-arithmetic-control}

\begin{itemize}
\item \op{11375} bit test on operands with high bits set, and field code 0's signedness above $2^{31}$ (\code{ledger/claims.toml}).
\item \op{10369}/10 mode-SET layout (operand 1 bit $2^{31}$): unreached from source (MM~\mm{25.73}, MM~\mm{25.76}).
\item Integer source-modifier value 8 as a signedness bit beyond \op{10239} and \op{16805} (MM~\mm{25.59}, MM~\mm{25.145.1}).
\item What \op{904} and \op{3818} do with a NaN input (MM~\mm{25.145.1}).
\item What the negate and abs modifiers do on \op{9700} (MM~\mm{25.145}).
\item How accurate \op{2570} is.
\item What the trig-family op \opn{3946} computes (glossary).
\item Whether \op{1048} is the fp32-to-bfloat narrowing (glossary).
\item Where the load-wait field of \op{798} sits (MM~\mm{25.196}).
\item How \op{9320} and \op{11179} wait, and how long they keep their sources (MM~\mm{25.196}).
\item Why a special-register read publishing on slot 7 stores nothing, whether or not it waits (MM~\mm{25.121}).
\item \code{SR\_PVSIMD} / \code{SR\_TVSIMD} under partial activity; roles of Apple's four-instruction divergent broadcast (MM~\mm{25.144.6}, MM~\mm{25.145.1}).
\item y and z builtins: no multi-axis dispatch has been run (MM~\mm{25.145.1}).
\item Exec auxiliary operand 0 meaning; count 3 and novel exec encodings: not executed (\code{isa/g17-exec-mask.txt}).
\item Whether \op{579} and \op{577} with n = 2 hold beyond the whole-program evidence they were admitted on (MM~\mm{25.145.2}).
\item How indirect calls execute; this compiler inlines every call (MM~\mm{25.90}).
\item Shuffles, broadcasts and votes in a partially active simdgroup (MM~\mm{25.145.1}); \op{14170} with a non-uniform mask (MM~\mm{12}).
\item Flag numeric value per producing form (\op{10379}, which has no glossary entry, against \op{11313}).
\end{itemize}

\section{Memory and synchronization}\label{sec:memory-synchronization-chapters}

\begin{itemize}
\item Out-of-range store behavior on hardware; other load forms and resource kinds out of range (MM~\mm{17} rule 34, \cref{sec:out-of-range-accesses-observed}).
\item Negative displacement; misaligned scalar access legality (\code{isa/g17-addressing-rules.json}).
\item Why Apple's compiler uses the index add, not the displacement, for unsigned \code{tid + k} (\code{isa/g17-addressing-rules.json}).
\item Displacement units of the \op{17256} (MM~\mm{25.141.18}).
\item Whether a masked-off load issues its access (no observable; \code{isa/g17-predication.json}).
\item Meaning of the threadgroup load's sensitive bits byte4{[}3{]}, byte4{[}4{]}, byte6{[}4{]} (\code{isa/g17-execution-tgbits-results.json}).
\item Hardware behavior for a threadgroup-memory declaration over 32 KiB, which Metal accepts (MM~\mm{25.147}).
\item Imageblock with more than one simdgroup or threadgroup; the store's operand-1 bit-4 flag; the read's 0x2000800000 word
  (MM~\mm{25.145.2}, MM~\mm{25.96}, MM~\mm{25.104.2}).
\item \op{12352} and \op{13276} address space (MM~\mm{25.144.1}).
\item Whether \op{17229} reads a texture fetch (Apple's tex2d-read); no dispatch has settled it (MM~\mm{25.146}).
\item Sampled and write texture classes executed (\code{agxforge/g17/texnarrow.py}).
\item Origin of \op{621}'s cap of four (\code{isa/g17-address-formula.txt}).
\item Cross-simdgroup winner of different-value same-address stores (MM~\mm{25.145.4}; \cref{ch:synchronization}).
\item Whether an execution-only barrier orders device memory by design (MM~\mm{14}).
\item Per-lane atomic min/max returns and unsigned long-form returns; atomics on non-rank-0 buffers; 64-bit atomics on hardware
  (\code{docs/g17-first-dispatch-trials.md}, \code{agxforge/g17/cc.py}).
\item Acquire/release semantics beyond the fence-count-fence pattern; emulator race check for atomics and \op{14156} (MM~\mm{25.197}).
\item Function of barrier-token variants 0x400 and 0x60000502 (MM~\mm{25.141.20}).
\item Overlap and ordering of multiple command records in one below-Metal Submit beyond the self-accumulation shape
  (\code{docs/g17-first-dispatch-trials.md}).
\item Barrier behavior in other partial threadgroup sizes and with inactive lanes (\code{docs/g17-first-dispatch-trials.md}).
\end{itemize}

\section{Tensor unit and performance}\label{sec:tensor-unit-performance}

\begin{itemize}
\item MMA-to-ALU result ordering: interlock or implicit wait (MM~\mm{6}, MM~\mm{25.90}).
\item What the destination keep flag does (no hardware effect found), and what population the 758,739 count covers (MM~\mm{0.4}, MM~\mm{25.22}).
\item Accumulator modifier (operand 10): a multi-bit encoding not excluded (MM~\mm{6}).
\item 36 narrow-accumulator declarations, \op{5108} through \op{5143}: refused although structurally like admitted forms (MM~\mm{25.4/F4} "F4, narrowed", MM~\mm{25.6}, MM~\mm{25.13}).
\item MM~\mm{11}'s "56 MMA forms" not reconciled with the 60 non-\code{\_shifted} declarations (MM~\mm{11}, MM~\mm{25.6}).
\item Bit-71 int8 MMA variant: executes bit-exact; field meaning and decoder identity unknown (MM~\mm{20}, MM~\mm{25.71}, row T11).
\item Eight decoder-refused bits of \op{5106} that change the result against six that do not: no encoding property separates them (MM~\mm{25.36}).
\item Whether out-of-place accumulation (D a different tuple from C) works; it has not been executed (MM~\mm{2}).
\item The int8 fragment byte order within the register pair (MM~\mm{0.7}).
\item Fragment shapes other than 16 × 16 × 16, which have not been measured (MM~\mm{0.14}).
\item Whether a release clears or frees per-lane storage (MM~\mm{25.89}, MM~\mm{25.95}).
\item What the legacy engine shares with fp32 fma ("8 fma per legacy issue" is a throughput equivalence) (MM~\mm{9}).
\item Software fp4/fp6 operand unpack: not built (MM~\mm{0.14}).
\item How the decode handles fp16 subnormals, which has not been measured (MM~\mm{0.14}).
\item Cause of the per-core simdgroup staircase inside the core (MM~\mm{25.122}).
\item Instruction-fetch cost at intermediate occupancy (between 1 and 26 simdgroups per core) (MM~\mm{25.144.12}).
\item Why Apple-compiled twins of this project's decode kernels are 10–15 percent faster (MM~\mm{25.211}).
\item How to identify per-simdgroup latency in the static model, and how to model operand sharing so that GEMM tile shapes extrapolate (MM~\mm{25.144.12}).
\item Whether the results hold on another G17 product or OS version; none has been measured (MM~\mm{0.1}, MM~\mm{20}).
\end{itemize}
